% !TeX root = ../papers/nc-p3.tex
\section{生成元の重みによる四つの分岐剰余類の正則性}
\label{sec:rank-fourteen-ramified-disks}
\begin{theorem}\label{thm:rank-fourteen-ramified-disks}
標数零の通常な階数14の窓は、$t$ の101進還元の一つが $19,80$ 以外なら
標準Koszul左右AS正則である。無限遠への還元も含む。従って未判定の有理閉点の
原始既約式 $P(t)$ の最高次係数は101で単元であり、
\[
 \overline{P}(t)=c(t-19)^a(t-80)^b,\qquad c\ne0,\quad a+b=\deg P
\]
が必要である。
\end{theorem}
\begin{proof}
前定理の四つの分岐基点 $d=-3,-1,1,\infty$ を扱う。
$101$ 進付値環を有限拡大して $t-d=w^{18}$、無限遠なら $t=w^{-18}$ とする。
四位置の各文字 $a$ に対し、テンソルのその一文字に作用する対角変換を
$w^{-e_{p,a}}$ と取る。次表の各欄は文字順の四つの重みである。
\[
\begin{array}{c|r|r|r|r}
d &p=0&p=1&p=2&p=3\\\hline
1&(-9,-9,1,-17)&(4,-14,12,12)&(2,2,-6,-24)&(-15,3,-7,-7)\\
-1&(-7,-7,3,-15)&(-3,-21,5,5)&(-3,-3,7,-11)&(-23,-5,3,3)\\
-3&(10,10,-7,11)&(-12,6,-13,-13)&(-1,-1,0,18)&(3,-15,2,2)\\
\infty&(8,8,-9,9)&(-12,6,-13,-13)&(9,9,-8,10)&(13,-5,-6,-6)
\end{array}
\]
分岐一様化変数 $u=w^9$ での各非零成分の位数 $v$ と先頭係数を前節と同じ式で取る。
\path{computations/verify_rank_fourteen_ramified_disks.py} は全成分で
$9v-\sum_p e_{p,a_p}\ge0$ を確認する。分母の先頭係数は101で単元なので、
この不等式は $w$ の正の付値の大きさによらず変換後の成分の整性を与える。
等号は各基点で32成分に限られ、いずれも $Y$ を含まない項が先頭である。
従って $Y$ の符号にも $w$ の選択にも依存せず、表から剰余テンソルを得る。

$i=10,91$ の各値で、変換後の剰余テンソルの一文字切断と隣接切断は
全て階数4と6を持つ。四始点で経路次元35,56、中心の文字次元 $(3,2,0,0)$、
四中心商の長さ13までの標準次元を確認する。分数体上では可逆な対角変換であり、
文字を保つため中心次元も保存する。一般正則点からの下界と還元上界による
定階数および中心持上げを適用すれば、元の分数体上の窓で全次数の主張を得る。
これと前定理で残る還元は19と80だけである。原始化した $P$ の斉次化で
最高次係数が消えれば無限遠に根を持ち、それも認証されるため、最後の必要条件が従う。
\end{proof}

\begin{proposition}
残る二剰余類では、上の文字基底の対角変換だけでは隣接階数6の還元を作れない。
六つの二次関係の平坦極限は長さ3,4,5で各始点の次元 $20,38,72$ を与える。
\end{proposition}
\begin{proof}
$t=-1+2i$ とその共役で $s=\pm1$ を代入する。四階テンソルの全64個の
文字許容成分はどれも101進単元で、他の成分は曲線上で恒等的に零である。
前節によりこの境界テンソルの隣接階数は4である。
同じ還元を持つ窓の各非零成分と対応する境界成分との比は付値環で1に還元される。
どの位置別対角変換を両者に施してもこの比は変わらず、共通の最小位数を払った後も
還元は一致する。定数境界テンソルの対角変換では階数は4以下にしかならない。

一方、各文字の右端ブロックを単位行列にした一般の関係式は、この二点の上で
全係数が正則である。そのまま特殊化した六行は独立で、関係六平面の極限を与える。
同じ検算コードで $\mathbb Q(i)$ 上、各文字の両符号と両共役について
これらの関係の全挿入を消去すると、主張の次元を得る。
\end{proof}
